\subsection{Proof of the flatness criterion}

\textbf{(A) The implications (i) \(\Rightarrow\) (ii) \(\Leftrightarrow\) (iii)}\par

The implication (i) \(\Rightarrow\) (ii) is immediate (Chapter I, § 4). The equivalence (ii) \(\Leftrightarrow\) (iii) is a special case of Chapter I, § 4, Proposition 2 applied to R = A, S = A/3, F = M, E = N, taking account of the fact that being given an (A/3)-module structure on N is equivalent to being given an A-module structure under which N is annihilated by 3.

Remark (1). Condition (ii) is also equivalent to the following:

(ii') \(\operatorname{Tor}_1^{\mathbf{A}}(\mathbf{N},\mathbf{M}) = 0\) for every A-module N annihilated by a power of \(\Im\) . Clearly (ii') implies (ii). Conversely, if (ii) holds, then in particular \(\operatorname{Tor}_1^{\mathbf{A}}(\Im^n\mathrm{N} / \Im^{n + 1}\mathrm{N},\mathbf{M}) = 0\) for all \(n\) ; from the exact sequence

\[
0 \rightarrow \Im^ {n + 1} N \rightarrow \Im^ {n} N \rightarrow \Im^ {n} N / \Im^ {n + 1} N \rightarrow 0
\]

we derive the exact sequence

\[
\operatorname{Tor} _ {1} ^ {\mathbf {A}} \left(\Im^ {n + 1} \mathrm{N}, \mathrm{M}\right)\rightarrow \operatorname{Tor} _ {1} ^ {\mathbf {A}} \left(\Im^ {n} \mathrm{N}, \mathrm{M}\right)\rightarrow \operatorname{Tor} _ {1} ^ {\mathbf {A}} \left(\Im^ {n} \mathrm{N} / \Im^ {n + 1} \mathrm{N}, \mathrm{M}\right)
\]

and, as there exists an integer \(m\) such that \(\mathfrak{S}^m\mathrm{N} = 0\) , we deduce by descending induction on \(n\) that \(\operatorname{Tor}_1^{\mathbf{A}}(\mathfrak{S}^n\mathrm{N},\mathbf{M}) = 0\) for all \(n\leqslant m\) and in particular for \(n = 0\) .

It follows from this that if \(\Im\) is nilpotent, (ii) implies (i), for (ii') then means that \(\operatorname{Tor}_1^{\mathbf{A}}(\mathbf{N},\mathbf{M}) = 0\) for every A-module \(\mathbf{N}\) and hence that \(\mathbf{M}\) is flat (Chapter I, §4).

\textbf{(B) Let us prove the following proposition:}\par

PROPOSITION 1. Let A be a commutative ring, \(\mathfrak{I}\) an ideal of A and M an A-module. The following conditions are equivalent:

\begin{enumerate}
\def\labelenumi{(\alph{enumi})}
\item
  For all \(n \geqslant 1\) , \(\operatorname{Tor}_1^{\mathbf{A}}(\mathbf{A}/\mathfrak{I}^n, \mathbf{M}) = 0\) .
\item
  For all \(n \geqslant 1\) , the canonical homomorphism
\end{enumerate}

\[
\theta_ {n} \colon \mathfrak {I} ^ {n} \otimes_ {\mathbf {A}} \mathbf {M} \rightarrow \mathfrak {I} ^ {n} \mathbf {M}
\]

is bijective.

Moreover these conditions imply:

\begin{enumerate}
\def\labelenumi{(\alph{enumi})}
\setcounter{enumi}{2}
\tightlist
\item
  The canonical homomorphism \(\gamma_{\mathrm{M}}: \operatorname{gr}(\mathbf{A}) \otimes_{\mathbf{g}_{0}(\mathbf{A})} \operatorname{gr}_{0}(\mathbf{M}) \to \operatorname{gr}(\mathbf{M})\) is bijective. Conversely, if \(\Im\) is nilpotent, (c) implies (a) and (b).
\end{enumerate}

The equivalence of (a) and (b) follows from the exact sequence

\[
0 = \operatorname{Tor} _ {1} ^ {\mathbf {A}} (\mathrm{A}, \mathrm{M}) \rightarrow \operatorname{Tor} _ {1} ^ {\mathbf {A}} (\mathrm{A} / \Im^ {n}, \mathrm{M}) \rightarrow \Im^ {n} \otimes_ {\mathbf {A}} \mathrm{M} \rightarrow \mathrm{M}.
\]

Consider next the diagram

\[
\begin{array}{c} \mathfrak {I} ^ {n + 1} \otimes_ {\mathrm{A}} \mathrm{M} \longrightarrow \mathfrak {I} ^ {n} \otimes_ {\mathrm{A}} \mathrm{M} \longrightarrow (\mathfrak {I} ^ {n} / \mathfrak {I} ^ {n + 1}) \otimes_ {\mathrm{A}} (\mathrm{M} / \mathfrak {I M}) \longrightarrow 0 \\ \theta_ {n + 1} \Bigg \downarrow \quad \theta_ {n} \Bigg \downarrow \quad \gamma_ {n} \Bigg \downarrow \\ 0 \longrightarrow \mathfrak {I} ^ {n + 1} \mathrm{M} \longrightarrow \mathfrak {I} ^ {n} \mathrm{M} \longrightarrow \operatorname{gr} _ {n} (\mathrm{M}) \longrightarrow 0 \end{array}\tag{1}
\]

where we note that \((\mathfrak{I}^n / \mathfrak{I}^{n+1}) \otimes_{\mathbf{A}} (\mathbf{M} / \mathfrak{I}\mathbf{M})\) is canonically identified with \((\mathfrak{I}^n / \mathfrak{I}^{n+1}) \otimes_{\mathbf{A} / \mathfrak{S}} (\mathbf{M} / \mathfrak{I}\mathbf{M})\) . This diagram is commutative by definition of \(\gamma_n\) and its rows are exact. If (b) holds, \(\theta_n\) and \(\theta_{n+1}\) are bijective and so therefore is \(\gamma_n\) by definition of cokernel, hence (b) implies (c). Conversely, assuming that \(\mathfrak{I}\) is nilpotent, let us show that (c) implies (b); we shall argue by descending induction on \(n\) , since \(\mathfrak{I}^n \otimes_{\mathbf{A}} \mathbf{M} = \mathfrak{I}^n\mathbf{M} = 0\) for \(n\) sufficiently large. Suppose then that in diagram (1), \(\gamma_n\) and \(\theta_{n+1}\) are bijective; then so is \(\theta_n\) by virtue of Chapter I, § 1, no. 4, Corollary 1 to Proposition 2.

\textbf{(C) The implication (ii) \(\Rightarrow\) (iv)}\par

If (ii) holds, so does (ii') by Remark 1; Proposition 1 then shows that \(\gamma_{\mathbf{M}}\) is an isomorphism. On the other hand, we already know that (ii) implies (iii) and hence M/Im is a flat (A/Im)-module, which completes the proof that (ii) implies (iv).

Remark (2). Proposition 1 shows that, if \(\mathfrak{S}\) is nilpotent, (iv) implies (iii); taking account of Remark 1, we have therefore proved in this case that (i), (ii), (iii) and (iv) are equivalent.

\textbf{(D) The equivalence (iv) \(\Leftrightarrow\) (v)}\par

For all \(n \geqslant 1\) , M/ \(\Im^{n}M\) has a canonical (A/ \(\Im^{n}\) )-module structure. If it is filtered by the \((\Im/\Im^{n})\) -adic filtration, it is immediate that \(\mathrm{gr}_{m}(\mathrm{M}/\Im^{n}\mathrm{M}) = \mathrm{gr}_{m}(\mathrm{M})\) if m \textless{} n and \(\mathrm{gr}_{m}(\mathrm{M}/\Im^{n}\mathrm{M}) = 0\) if \(m \geqslant n\) . For all \(k \geqslant 1\) , let \(A_{k} = A/\Im^{k}\) , \(\Im_{k} = \Im/\Im^{k} M_{k} = M/\Im^{k} M\) ; let \((\mathrm{iv})_{k}\) (resp. \((\mathrm{v})_{k}\) ) denote the assertion derived from (iv) (resp. (v)) by replacing A, \(\Im\) , M by \(A_{k}\) , \(\Im_{k}\) , \(M_{k}\) . It follows from what has just been said that (iv) is equivalent to ``for all \(k \geqslant 1\) , \((\mathrm{iv})_{k}\) '' and obviously (v) is equivalent to ``for all \(k \geqslant 1\) , \((\mathrm{v})_{k}\) ''. Then it will suffice to establish the equivalence \((\mathrm{iv})_{k} \Leftrightarrow (\mathrm{v})_{k}\) for all k or also to show that (iv) \(\Leftrightarrow (\mathrm{v})\) when \(\Im\) is nilpotent. Now (Remark 2) we have seen that in that case (iv) is equivalent to (i). As M/ \(\Im^{n}M\) is isomorphic to M \(\otimes_{A} (A/\Im^{n})\) , (i) implies (v) (Chapter I, §2, no. 7, Corollary 2 to Proposition 8); moreover clearly (v) then implies (i). We have therefore shown the equivalence (iv) \(\Leftrightarrow (\mathrm{v})\) in all cases and also that of all the properties of the theorem in the case where \(\Im\) is nilpotent.

\textbf{(E) The implication \((\mathbf{v})\Rightarrow (\mathbf{i})\) when \(\mathbf{A}\) is Noetherian and \(\mathbf{M}\) ideally Hausdorff}\par

It is sufficient to prove that for every ideal \(\mathfrak{a}\) of \(\mathbf{A}\) the canonical mapping \(j\colon \mathfrak{a} \otimes_{\mathbf{A}} \mathbf{M} \to \mathbf{M}\) is injective (Chapter I, § 2, no. 3, Proposition 1). Let \(x \in \operatorname{Ker} j\) ; as \(\mathfrak{a} \otimes_{\mathbf{A}} \mathbf{M}\) is Hausdorff with the \(\mathfrak{S}\) -adic topology, it suffices to verify that, for every integer \(n > 0\) , \(x \in \mathfrak{S}^n(\mathfrak{a} \otimes_{\mathbf{A}} \mathbf{M})\) . Let \(f\colon \mathfrak{S}^n\mathfrak{a} \to \mathfrak{a}\) be the canonical injection; it suffices to show that \(x \in \operatorname{Im}(f \otimes 1_{\mathbf{M}})\) ; for if \(b \in \mathfrak{S}^n\) , \(a \in \mathfrak{a}\) and \(m \in \mathbf{M}\) , the image under \(f \otimes 1_{\mathbf{M}}\) of the element \((ba) \otimes m\) of \((\mathfrak{S}^n\mathfrak{a}) \otimes_{\mathbf{A}} \mathbf{M}\) is the element \((ba) \otimes m = b(a \otimes m)\) of \(\mathfrak{a} \otimes_{\mathbf{A}} \mathbf{M}\) and hence \(\operatorname{Im}(f \otimes 1_{\mathbf{M}}) \subset \mathfrak{S}^n(\mathfrak{a} \otimes_{\mathbf{A}} \mathbf{M})\) . By virtue of Krull's Theorem (§ 3, no. 2, Theorem 2), there exists an integer \(k\) such that \(\mathfrak{a}_k = \mathfrak{a} \cap \mathfrak{S}^k \subset \mathfrak{S}^n\mathfrak{a}\) ; if \(i: \mathfrak{a}_k \to \mathfrak{a}\) is the canonical injection, it will then be sufficient to show that \(x \in \operatorname{Im}(i \otimes 1_M)\) . Now, denoting by \(p: \mathfrak{a} \to \mathfrak{a}/\mathfrak{a}_k\) and \(h: \mathfrak{a}/\mathfrak{a}_k \to \mathrm{A}/\mathfrak{S}^k\) the canonical mappings, there is a commutative diagram

\[
\begin{array}{c} \mathfrak {a} _ {k} \otimes_ {\mathrm{A}} \mathrm{M} \xrightarrow {i \otimes 1 _ {\mathrm{M}}} \mathfrak {a} \otimes_ {\mathrm{A}} \mathrm{M} \xrightarrow {p \otimes 1 _ {\mathrm{M}}} (\mathfrak {a} / \mathfrak {a} _ {k}) \otimes_ {\mathrm{A}} \mathrm{M} \longrightarrow 0 \\ j \Biggl \downarrow \\ \mathrm{M} \quad \longrightarrow \quad (\mathrm{A} / \mathfrak {I} ^ {k}) \otimes_ {\mathrm{A}} \mathrm{M} \end{array}
\]

in which the first row is exact. It suffices to prove that \(x \in \operatorname{Ker}(p \otimes 1_{\mathbf{M}})\) and, as \(x \in \operatorname{Ker} j\) by hypothesis, it will suffice to verify that the mapping \(h \otimes 1_{\mathbf{M}}\) is injective. Now, it may also be written (Algebra, Chapter II, § 3, no. 6, Corollary 3 to Proposition 6)

\[
h \otimes 1 _ {\mathbf {M} / \Im^ {k} \mathbf {M}}: (a / a _ {k}) \otimes_ {\mathbf {A} / \Im^ {k}} (\mathbf {M} / \Im^ {k} \mathbf {M}) \rightarrow \mathbf {M} / \Im^ {k} \mathbf {M}
\]

and, as \(h\) is injective and, by (v), \(\mathbf{M} / \Im^k\mathbf{M}\) is a flat \((\mathbf{A} / \Im^k)\) -module, this completes the proof.

